\chapter{Objetivos}
\label{cap:objetivos}

Los objetivos de esta investigación se derivan directamente de la pregunta formulada en el planteamiento del problema (capítulo~\ref{cap:problema}) y del vacío de conocimiento identificado en los antecedentes: la ausencia de evidencia experimental consistente y reproducible sobre el consumo energético de distintas formas de organizar un sistema de software. En coherencia con el problema planteado, los objetivos no buscan demostrar la superioridad de una arquitectura sobre otra, sino establecer si existen diferencias entre ellas.

\section{Objetivo general}
\label{sec:objetivo-general}

Determinar si existen diferencias en el consumo energético del software entre una arquitectura monolítica tradicional y una arquitectura de monolito modular, bajo condiciones y cargas de trabajo similares.

\section{Objetivos específicos}
\label{sec:objetivos-especificos}

\begin{enumerate}
    \item Caracterizar la arquitectura de monolito tradicional y la de monolito modular, estableciendo las condiciones de equivalencia funcional, de carga de trabajo, de hardware y de entorno de ejecución bajo las cuales serán comparadas.
    \item Medir el comportamiento energético de cada arquitectura ---utilización de CPU, uso de memoria, tiempo de ejecución y consumo energético por unidad de trabajo--- bajo condiciones controladas y comparables.
    \item Comparar los resultados obtenidos en ambas arquitecturas para establecer si el tipo de arquitectura monolítica se asocia con diferencias en el consumo energético.
    \item Analizar las implicaciones de los resultados para la sostenibilidad ambiental del software y su posible uso como criterio en la toma de decisiones sobre el diseño de software.
\end{enumerate}